\documentclass[11pt]{article}

\usepackage[margin=1in]{geometry}
\usepackage{amsmath,amssymb,amsthm}
\usepackage{booktabs}
\usepackage{tabularx}
\usepackage{hyperref}
\usepackage{cleveref}
\usepackage{enumitem}
\usepackage[utf8]{inputenc}
\usepackage[T1]{fontenc}

\newtheorem{theorem}{Theorem}

\usepackage[
    type={CC},
    modifier={by-nc},
    version={3.0},
]{doclicense}

\title{Don't Let Your System Decide It's Dead}
\author{Aaron H. Alpar\\Independent Researcher}
\date{}

\begin{document}
\doclicenseThis
\maketitle

\begin{abstract}
Distributed systems fail in two ways: the failure itself, and the system's
automatic response to the failure. The second is often worse. The root cause is
systems making semantic decisions with insufficient information---promoting
``slow'' to ``dead'' based on wall-clock constants rather than measured resource
models. This paper identifies a precise invariant---the \emph{recovery
invariant}---that separates safe automatic recovery from cascading failure, and
a design principle---the \emph{compensation boundary}---that determines who
should make which decisions. The recovery invariant has three levels: individual
cost boundedness, aggregate rate boundedness, and feedback stability (spectral
radius of the recovery gain matrix less than one). The compensation boundary is
the architectural point where a system stops surfacing failure to a higher
authority and starts acting on it automatically. We analyze TCP's retransmission
timeout as a reference design that satisfies all three levels, and show how
Cassandra's hinted handoff, UNIX swap and OOM killer, and wall-clock-based
tombstone garbage collection violate them. We connect this framework to the
metastable failures literature, resilience engineering, and control theory, and
apply the principles to CRDT tombstone garbage collection as a case study.
\end{abstract}

\section{Introduction}

Distributed systems fail in two ways: the failure itself, and the system's
automatic response to the failure. The second is often worse.

The root cause is systems making semantic decisions with insufficient
information---promoting ``slow'' to ``dead'' based on wall-clock constants
rather than measured resource models. There is a precise invariant that
separates safe automatic recovery from cascading failure, and a design principle
that tells you who should make which decisions. Most production systems violate
both.

\section{The Setup: A Configuration Parameter That Eats Your Data}

Cassandra has a setting called \texttt{gc\_grace\_seconds}. It defaults to
864,000---ten days. Here is what it does: when you delete a row, Cassandra does
not remove it. It writes a \emph{tombstone}---a marker that says ``this key was
deleted.'' Tombstones exist because in an eventually consistent system, if you
just remove the data, another replica that has not heard about the deletion will
happily replicate the old value back. The tombstone is the deletion's proof of
existence.

\texttt{gc\_grace\_seconds} is how long Cassandra keeps that proof around.
After ten days, Cassandra garbage-collects the tombstone. The assumption: repair
has propagated the tombstone to all replicas within that window.

If repair did not run---or did not finish---within those ten days, the tombstone
disappears from some replicas while the pre-delete value still exists on others.
The deleted data silently reappears. This is a well-documented production
failure mode.

What just happened? A single wall-clock constant made two different decisions
simultaneously:

\begin{enumerate}
  \item \textbf{A local decision}: ``This tombstone's storage can be
    reclaimed.'' (Garbage collection---the system has the information to do
    this.)
  \item \textbf{A semantic decision}: ``All replicas have observed this
    deletion.'' (Replication state---the system is \emph{guessing}.)
\end{enumerate}

These are different decisions requiring different information. Conflating them is
the root of most cascading failures in distributed systems, and the pattern
repeats everywhere: timeouts that promote ``slow'' to ``dead,'' OOM killers that
choose which process to sacrifice, swap systems that silently trade memory
bandwidth for disk bandwidth until everything grinds to a halt.

This paper is about drawing the line correctly.

\section{The Two Decisions}

Every failure in a distributed system requires two decisions:

\begin{enumerate}
  \item \textbf{The local decision}: ``This connection/request/operation has
    failed.''
  \item \textbf{The semantic decision}: ``This node/peer/service is gone and I
    should act accordingly.''
\end{enumerate}

The first is a statement about what the system has observed. The second is an
interpretation of what that observation \emph{means}. They require different
information, are made at different layers, and have different consequences when
wrong.

Getting the local decision wrong is usually cheap---a spurious timeout, a wasted
retry. Getting the semantic decision wrong can be catastrophic---deleted data
resurrected, a healthy node declared dead and its traffic redistributed into an
already-overloaded cluster, a cascade that turns a minor hiccup into a major
outage.

The engineering challenge is that the semantic decision is \emph{tempting} to
automate. Manual intervention is slow. Operators are expensive. And most of the
time, the guess is right---the node really is dead, repair really did finish,
the process really was misbehaving. Systems are designed for the common case,
and the common case makes automation look safe. The failure modes live in the gap
between ``usually right'' and ``always right.''

\section{TCP Got This Right (Mostly)}

TCP's retransmission timeout is the gold standard for automatic failure
decisions. It gets nearly everything right, and examining \emph{why} reveals the
principles that other systems violate.

\textbf{Derived from measurement, not magic constants.} TCP's retransmission
timeout (RTO) is computed from Jacobson's round-trip time (RTT) estimator---a
smoothed average plus variance of measured round-trip times. The timeout tracks
actual network behavior, not an engineer's guess about what it should be.

\textbf{Individual recovery cost is bounded and known.} A single retransmission
costs one segment---fixed, known in advance, expressed in the same units as
normal operation.

\textbf{Recovery never amplifies the failure.} Exponential backoff ensures the
retry \emph{rate} converges to zero. Each successive timeout doubles the wait.
The more congested the network, the less recovery traffic TCP adds to it.

\textbf{Failure is surfaced, not hidden.} When TCP exhausts its retries, the
application gets \texttt{ETIMEDOUT}. TCP does not guess what the timeout
\emph{means}---it tells the application ``this connection is dead'' and lets the
application decide whether to reconnect, try a different server, or alert an
operator.

TCP separates the layers cleanly:

\begin{itemize}
  \item \textbf{TCP decides}: ``This connection is dead.'' (It has the
    information: measured RTT, retry history, bounded retransmission cost.)
  \item \textbf{Application decides}: ``What does this dead connection mean?''
    (It has the context: service topology, redundancy, retry policy.)
  \item \textbf{Neither decides}: ``This node is permanently gone.'' (This
    requires physical-world knowledge---hardware state, operator intent---that
    no software layer has.)
\end{itemize}

There is a nuance most people miss: TCP \emph{does} eventually give up
automatically. \texttt{tcp\_retries2} defaults to 15 retries, which works out to
roughly 13--30 minutes depending on measured RTT. This is a wall-clock timeout,
but grounded in a measured resource model---derived from observed network
behavior, not pulled from thin air.

\subsection{Why This Separation Is Provably Necessary}
\label{sec:flp}

\subsubsection{The Impossibility Result}

The FLP impossibility result~\cite{fischer1985impossibility} proves that in an
asynchronous system where at least one process may crash, no deterministic
algorithm can guarantee consensus. The load-bearing assumption is
\emph{asynchrony}: with no bound on message delay or processing time, you cannot
distinguish a crashed process from a slow one.

TCP operates \emph{within} the bounds where this distinction is tractable---it
measures RTT, so it has a probabilistic model of ``how slow is normal.'' But the
question ``is this node permanently gone?'' lives in the asynchronous regime
where FLP applies. No amount of measurement resolves it, because the relevant
information (hardware failure, network partition, operator intent) is outside the
system's observation boundary.

This is the theoretical foundation for the rest of the argument: systems that
promote ``slow'' to ``dead'' automatically are not just making a bad engineering
choice---they are making a decision that is \emph{provably underdetermined} by
the information available to them. The correct response is the one TCP models:
decide what you can (connection liveness), surface what you cannot
(\texttt{ETIMEDOUT}), and let the layer with more information make the semantic
call.

\subsubsection{The Workaround}

Chandra and Toueg~\cite{chandra1996unreliable} formalized the workaround: you
can solve consensus if you have a \emph{failure detector} with known properties
(completeness and accuracy). Different detector strengths enable different
guarantees. This maps directly to the layered model above---each layer is a
failure detector with a specific information boundary, and the system's
correctness depends on no layer exceeding its boundary.

\section{The Recovery Invariant}

TCP's design is not just good engineering intuition---it satisfies a precise
invariant that separates safe automatic recovery from cascading failure. The
invariant has three levels. Each is necessary. Together they guarantee bounded
recovery in the linearized model developed in \cref{sec:gain-matrix};
\cref{sec:nonlinearity} bounds the gap to the nonlinear system.

\subsection{Precondition: You Need a Resource Model}

Before evaluating any level, the system must answer a basic question:
\emph{recovery costs what, measured in what units, against what capacity?}

The system must name its bottleneck resource, measure its total capacity, and
express both normal operations and recovery actions in the same units against
that capacity.

TCP has an explicit resource model: bandwidth, measured in segments per
round-trip. One normal packet and one retransmission cost the same unit---one
segment---against a known link capacity.

Cassandra's hinted handoff has a partial model: one hint replay costs roughly
one write in I/O operations per second (IOPS). But the available IOPS during
replay depend on concurrent read load, compaction, and repair---variables the
hint replay system does not observe.

The OOM killer has almost no resource model: it knows system-wide memory
pressure triggered the kill, but the cost of killing a specific process (how
much memory is freed, how long until the supervisor restarts it, what the restart
will allocate) depends on external state the kernel does not track.

The precondition is what separates systems that \emph{can} reason about recovery
cost from those that act and hope.

\subsubsection{Formal Framework: Bandwidth as Primitive}
\label{sec:bandwidth}

Every system resource has a consumption rate expressible as bandwidth:
\textbf{units of resource per unit time}.

\begin{center}
\begin{tabular}{@{}ll@{}}
\toprule
Resource & Bandwidth unit \\
\midrule
Network  & bytes/sec or segments/RTT \\
Disk I/O & IOPS or bytes/sec \\
CPU      & utilization (dimensionless rate) \\
Memory   & pages/sec (fault-in rate) \\
\bottomrule
\end{tabular}
\end{center}

For pure-flow resources (network, CPU), bandwidth \emph{is} the resource. There
is no at-rest state---a CPU cycle is consumed or not, a packet is in flight or
not.

For stock resources (memory capacity, disk space), the at-rest state is the
integral of past bandwidth:
\begin{equation}
  S(t) = \int_0^t b_{\mathrm{net}}(\tau)\,d\tau
\end{equation}
where $b_{\mathrm{net}}$ is the net consumption rate (allocation rate minus free
rate). Every physical memory page was populated through a page fault---zero-fill
or disk read---that consumed memory bus bandwidth and potentially disk bandwidth.
Virtual address space can be reserved without bandwidth (\texttt{malloc} with
overcommit returns immediately), but physical resource consumption always goes
through bandwidth.

This unification matters: the recovery invariant constrains rates. Stocks enter
only as boundary conditions on accumulated rates.

\subsubsection{Multi-Resource Model}
\label{sec:multi-resource}

A system has $n$ resources. At time $t$:

\begin{itemize}[nosep]
  \item $C_i$: capacity of resource $i$ (bandwidth ceiling, in
    $\mathrm{units}_i/\mathrm{sec}$)
  \item $w_i(t)$: bandwidth consumed by normal operations on resource $i$
  \item $r_i(t)$: bandwidth consumed by recovery actions on resource $i$
  \item $h_i(t) = C_i - w_i(t) - r_i(t)$: headroom on resource $i$
\end{itemize}

These form vectors $\mathbf{C}$, $\mathbf{w}(t)$, $\mathbf{r}(t)$,
$\mathbf{h}(t) \in \mathbb{R}^n$. For stock resources, an additional integral
constraint applies:
\begin{equation}
  \int_0^t b_{i,\mathrm{net}}(\tau)\,d\tau \leq S_i
\end{equation}

Recovery is safe when both conditions hold on every resource:
\begin{enumerate}
  \item $h_i(t) \geq 0$ at all $t$ (bandwidth not exceeded)
  \item $\int_0^t b_{i,\mathrm{net}}(\tau)\,d\tau \leq S_i$ for stock resources
    (capacity not exhausted)
\end{enumerate}

The bottleneck is whichever resource hits either constraint first. It is
emergent, not fixed---recovery actions can shift the bottleneck from one resource
to another.

\subsection{Level 1: Individual Bound}

\begin{quote}
Each recovery action's resource cost must be bounded and known.
\end{quote}

This is the weakest condition. You must know what a single compensating action
costs before you trigger it.

TCP satisfies this: one retransmission costs one segment---fixed, known,
independent of system state. Cassandra's hinted handoff satisfies this too: one
hint replay costs approximately one write. The OOM killer satisfies this: one
kill frees the memory of one process.

Every system discussed in this paper passes level~1. That is why it is
insufficient. A system can have individually cheap recovery actions and still
cascade---the danger is in their aggregate behavior.

\subsection{Level 2: Aggregate Bound}

\begin{quote}
At any point in time, the aggregate resource consumption by all concurrent
recovery actions must remain within the headroom not consumed by normal
operations: $h_i(t) = C_i - w_i(t) - r_i(t) \geq 0$ for every resource~$i$.
\end{quote}

This is the burst dimension. The distinction between ``amortized OK'' and
``worst-case catastrophic'' lives here.

TCP satisfies this: exponential backoff ensures the \emph{rate} of
retransmissions decreases with each round, so aggregate retry bandwidth shrinks
over time. Google's SRE retry budget is an explicit level-2 mechanism---capping
retries at 10\% of total requests bounds the aggregate cost to 1.1$\times$
normal load regardless of failure count.

Cassandra's hinted handoff violates this: hints accumulate linearly during a
node's absence and replay in a burst on return. A 10-minute outage at 100~Mbps
per node produces ${\sim}7$~GiB of hints; at the default 1024~KB/s throttle,
replay takes ${\sim}2$~hours. If the throttle is raised (as operators under
pressure often do), the burst can recreate the overload that caused the original
failure.

\subsection{Level 3: Feedback Bound}

\begin{quote}
Recovery actions must not create conditions that trigger further failures.
Equivalently: the recovery feedback loop must have gain $< 1$.
\end{quote}

Levels 1 and 2 treat recovery actions as independent events. Level~3 asks
whether they interact. A system violates level~3 when recovery actions consume
the same resources that are under contention, creating a positive feedback loop:
failure $\to$ recovery $\to$ more contention $\to$ more failure.

TCP satisfies this: backoff \emph{reduces} congestion on the link, which reduces
the condition (congestion) that caused the timeout in the first place. The
feedback loop is damped---each iteration moves the system closer to equilibrium,
not further.

The OOM killer can violate this: a process supervisor restarts the killed
process, which allocates memory, which recreates the pressure that triggered the
kill. The feedback loop gain depends on external configuration (supervisor
policy) that the kernel does not know about---another instance of a system making
decisions outside its information boundary.

Cassandra's hinted handoff can violate this: burst replay causes CPU/IO
contention $\to$ read latencies spike $\to$ read timeouts trigger read repairs
$\to$ read repairs consume more CPU/IO $\to$ more timeouts. The recovery action
for writes triggers a cascade in the read path.

\paragraph{The decompensation cascade} is what a level~3 violation looks like in
practice:

\begin{enumerate}
  \item System overloaded $\to$ operations slow down
  \item Timeouts fire $\to$ compensating actions triggered
  \item Compensating actions consume the bottleneck resource $\to$ more
    contention
  \item More timeouts $\to$ more compensating actions $\to$ feedback gain
    $\geq 1$ $\to$ cascade
\end{enumerate}

Systems designed for ``graceful degradation'' often degrade ungracefully because
the compensating actions were designed for the steady-state cost model (level~1
holds, level~2 holds at low utilization), not the failure cost model (level~2
breaks under burst, level~3 breaks under sustained pressure). David
Woods~\cite{woods2015four,woods2018theory} calls this \emph{decompensation}---the
exhaustion of adaptive capacity when the system's own recovery mechanisms become
the dominant source of load.

The metastable failures literature~\cite{bronson2021metastable,huang2022metastable}
identifies the same phenomenon: ``the root cause of a metastable failure is the
sustaining feedback loop, rather than the trigger.'' Their taxonomy of triggers
(load-spiking vs.\ capacity-decreasing) and amplification mechanisms (workload
amplification vs.\ capacity degradation) maps directly onto levels~2 and~3 of
this invariant. The formalization in \cref{sec:gain-matrix} approaches the same
problem from a different mathematical tradition---control theory and linear
algebra rather than Markov chains and queuing theory.

\paragraph{A note on quiesced recovery.} Some systems sidestep levels~2 and~3
entirely by deferring recovery to a state where no production traffic competes
for resources---replaying a write-ahead log on startup before accepting
connections, running repair during a maintenance window, or batching
reconciliation while the node is offline. When the competing workload is zero,
level~2 is trivially satisfied and level~3's feedback loop is broken. This is
effectively admission control: the system refuses new work until recovery is
complete. The trade-off is availability---the node is down during recovery.

\subsection{Summary}

\begin{small}
\noindent\begin{tabularx}{\textwidth}{@{}lXXXX@{}}
\toprule
Level & Condition & TCP & Hinted handoff & OOM killer \\
\midrule
0 & Resource model: $C_i$, $w_i(t)$, $r_i(t)$ in shared bandwidth units
  & \checkmark\ segments on a link
  & $\times$\ partial: IOPS but not concurrent load
  & $\times$\ memory freed unknown until kill \\
1 & Individual: cost of one action bounded and known
  & \checkmark\ one segment
  & \checkmark\ one write
  & \checkmark\ one kill \\
2 & Aggregate: recovery rate $\leq$ freed capacity
  & \checkmark\ exponential backoff
  & $\times$\ burst on return
  & \checkmark\ single kill \\
3 & Feedback: recovery gain $< 1$
  & \checkmark\ backoff reduces congestion
  & $\times$\ replay triggers read repair cascade
  & $\times$\ supervisor restarts killed process \\
\bottomrule
\end{tabularx}
\end{small}

The recovery invariant tells you whether recovery \emph{is} safe. It does not
tell you whether it will \emph{stay} safe, or what to do when it is not. A
system can satisfy all three levels at current load and violate them minutes
later. That question---when to compensate automatically and when to stop and
escalate---requires a different concept.

\section{The Compensation Boundary}

\begin{quote}
The \emph{compensation boundary}\footnote{The term \emph{compensating action}
originates in transaction processing, where a compensating transaction undoes
the effects of a committed transaction that cannot be rolled
back~\cite{gray1993transaction}. \emph{Decompensation}---the exhaustion of
adaptive capacity when recovery mechanisms become the dominant load source---is
David Woods' contribution (see \cref{sec:related-work}). The \emph{compensation
boundary} as used here---the architectural point where the system stops
surfacing failure and starts acting on it---is a framing specific to this paper,
combining the transactional concept (what is a compensating action?) with the
resilience engineering concept (when does compensation become the problem?) and
adding the design question: \emph{where should the system draw that line?}} is
the point in a system's architecture where it stops surfacing failure to a
higher authority and starts acting on it automatically. Below the boundary,
errors are reported---the caller decides. Above it, the system compensates
silently---no one is asked.
\end{quote}

The placement of that boundary determines how the system fails under pressure.

\subsection{z/OS: The Boundary Drawn Furthest Out}

z/OS draws the compensation boundary furthest out of any system I know. WLM
assigns work to service classes with defined goals; when the system cannot meet
goals, it reports the shortfall rather than silently
degrading.\footnote{WLM (Workload Manager) is the z/OS component that classifies
work into service classes with defined performance goals (response time,
velocity, discretionary) and dynamically manages system resources---CPU
dispatching priority, I/O priority, memory---to meet those goals. When goals
cannot be met, WLM reports the shortfall rather than silently degrading. See
IBM, \emph{z/OS MVS Planning: Workload Management}, SC34-2662.} z/OS manages
memory through explicit region sizes---exceed yours and the job abends (S878,
S80A), surfacing a hard failure to the operator.

XCF heartbeat failure detection declares a member dead (it has the information:
heartbeat history, configurable intervals), but does not decide what the failure
\emph{means} for application-level data---recovery exits handle that decision,
and may themselves escalate.\footnote{XCF (Cross-system Coupling Facility)
provides group membership, signaling, and heartbeat-based failure detection
across z/OS systems within a sysplex. XCF monitors member health via
configurable heartbeat intervals and declares a member failed when heartbeats
stop---a local decision about connectivity. What that failure \emph{means} for
application-level state is delegated to recovery exits registered by the
application, not decided by XCF itself. See IBM, \emph{z/OS MVS Setting Up a
Sysplex}, SA23-1399.} Automation is policy-driven (SA z/OS): an operator
authored the rule, every automated action is logged, and when no policy matches,
the system escalates rather than guessing.\footnote{SA z/OS (IBM System
Automation for z/OS) is a policy-driven automation framework for z/OS
operations. Operators author automation rules that define responses to system
events; every automated action is logged and auditable. When no rule matches a
condition, SA z/OS escalates to the operator rather than guessing. See IBM,
\emph{System Automation for z/OS Planning and Installation}, SC34-2571.}

The escalation mechanism itself is significant: WTO/WTOR (Write to Operator with
Reply) is \emph{synchronous}---on critical decisions, the system issues a
message and \textbf{waits for the operator to respond} before proceeding. The
operator is in the control loop by design---consulted during the decision, not
notified afterward.

But z/OS is not immune. Under system-wide memory pressure, it pages to auxiliary
storage silently. SRM---which coexists with WLM as its resource-allocation
mechanism---swaps out entire address spaces based on system-wide monitoring
(paging rates, UIC values, central storage demand). WLM sets policy goals; SRM
decides the mechanics. That separation is itself an example of decision authority
tracking information: WLM knows what matters to the workload, SRM knows what is
happening to the hardware. But the swap-out is still an automatic decision about
who gets to run. The compensation boundary is further out than UNIX's, not
absent.

\subsection{UNIX: A System at War with Itself}

UNIX draws the boundary closer in, and the result is a system with a split
personality.

The syscall interface surfaces errors faithfully---\texttt{errno},
\texttt{ENOSPC}, \texttt{ENOMEM} from \texttt{fork}, \texttt{SIGPIPE}, exit
codes. This is the ``let the caller decide'' philosophy. But the resource
management layer behind it hides scarcity, and virtual memory is the clearest
example.

When physical memory is exhausted, UNIX pages to disk---a compensating action
designed to degrade rather than fail. The capacity model is straightforward:
physical RAM plus swap, both measurable in bytes, upper bound known at boot. But
the \emph{cost} model is incoherent. The system is not trading memory for disk
\emph{space} (swap is pre-allocated); it is trading memory bandwidth for disk
bandwidth. A memory access costs ${\sim}100$\,ns; a page fault to an HDD costs
${\sim}10$\,ms; to NVMe, ${\sim}100$\,\textmu s. The compensating action
substitutes a resource that is 1,000--100,000$\times$ slower per access, and the
only unit that makes both sides commensurable is time---the one unit the system
does not budget.

Against the recovery invariant: Level~1 holds---one page swap frees one page of
physical memory, bounded and known. Level~2 is conditional---it holds while disk
I/O remains below saturation, but the system swaps without checking disk
bandwidth. Level~3 breaks under pressure: swapping consumes both memory bandwidth
(page copies) and disk bandwidth (I/O operations), slowing everything $\to$
queues build $\to$ more memory allocated to pending work $\to$ more swapping.
Thrashing is a level~3 violation---the recovery action amplifies the condition it
compensated for.

This is a defensible decision within UNIX's information boundary---the kernel
knows about memory and disk, and paging is within its operational scope. It makes
no semantic guess about the application. But acting within your boundary does not
guarantee stability. Swap satisfies the invariant at low memory pressure and
violates it under high pressure, and the system has no mechanism to predict the
transition. The implicit stabilization mechanism is degradation itself: as
everything slows, users give up, load drops. This is the \emph{absence} of
admission control, not a form of it.

The OOM killer is where the contradiction becomes explicit. When oversubscription
hits the wall, Linux makes an automatic semantic decision with unbounded
consequences---it kills a \emph{different} process than the one that caused the
pressure, chosen by heuristic (\texttt{oom\_badness} scoring, primarily by
memory footprint). Even the default overcommit mode (mode~0, heuristic) allows
this; mode~1 (\texttt{vm.overcommit\_memory=1}) is starker: \texttt{malloc}
\emph{always} succeeds regardless of available memory, then the kernel kills
processes later when it cannot honor the promise. Kubernetes sets mode~1 by
default.

The error is not surfaced---it is deferred until it is unrecoverable. When the
OOM killer fires, it logs the kill to the kernel ring buffer and
continues---no operator acknowledgment, no pause before acting.

UNIX does have tools that push the boundary outward: disk quotas
(\texttt{EDQUOT}), ulimits (\texttt{SIGXCPU}, \texttt{SIGXFSZ},
\texttt{EMFILE}), cgroups (scoped OOM, CPU throttling). When configured, these
behave like the mainframe model---the system says ``no'' and surfaces an error.
But they are all opt-in. A fresh UNIX system has no disk quotas, generous
ulimits, no cgroups. Compare z/OS: you cannot submit a job without a region
size; WLM service classes are mandatory. The compensation boundary in UNIX is a
configuration choice. In z/OS it is a design invariant---the system will not run
without it. Under production pressure, opt-in boundaries tend to be the first
thing relaxed or forgotten.

\subsection{Dynamo-Family Systems: No Boundary at All}

Dynamo-family systems (Cassandra, Riak, early DynamoDB) barely have a
compensation boundary. Compensating actions are implicit in each subsystem, not
centrally authored or auditable. Cassandra's \texttt{gc\_grace\_seconds} promotes
``old'' to ``replicated everywhere'' based on a wall-clock constant. Hinted
handoff accumulates hints during a node failure and replays them in a burst on
return. No escalation path exists---when the automatic decision is wrong, the
system has already acted.

\subsection{The Pattern}

\begin{small}
\noindent\begin{tabularx}{\textwidth}{@{}lXXXX@{}}
\toprule
 & Reports failure & Compensates silently & Escalation model & Structural? \\
\midrule
z/OS & Abends, WLM goal misses, console messages & Paging under system-wide pressure
  & Synchronous---WTO/WTOR waits for operator & Yes---cannot run without region sizes, WLM \\
UNIX & \texttt{errno} at syscall boundary & VM overcommit, swap, OOM killer
  & Asynchronous---syslog/dmesg, after the fact & No---quotas/cgroups exist but are opt-in \\
Dynamo & Logs (if watching) & Everywhere: GC, hinted handoff, read repair
  & None---no escalation path & No---compensation is the design \\
\bottomrule
\end{tabularx}
\end{small}

As the compensation boundary moves inward, the system's failure modes shift from
\emph{recoverable} (operator gets a message, decides what to do) to
\emph{unrecoverable} (data silently lost, wrong process killed, cascade already
in progress). The key property is not whether a system automates---z/OS automates
extensively---but whether \emph{wrong} automatic decisions are surfaced before
their consequences become permanent.

\section{Decision Authority Tracks Information Asymmetry}

The design principle that ties the preceding sections together:

\begin{quote}
Automatic decisions are safe when the decider has sufficient information to
guarantee the recovery invariant. When information is insufficient, the system
must surface state to the entity that \emph{can} distinguish---and wait.
\end{quote}

\begin{small}
\noindent\begin{tabularx}{\textwidth}{@{}lXl@{}}
\toprule
Decider & Has information about & Can safely decide \\
\midrule
TCP & Measured RTT, retry history, bounded retransmission cost
  & ``Connection is dead'' \\
Application & Service topology, redundancy, retry policy
  & ``Retry elsewhere'' \\
Operator & Physical world, business context, acceptable risk
  & ``Node is permanently gone'' \\
\bottomrule
\end{tabularx}
\end{small}

Each layer can make decisions within its information boundary. No layer should
make decisions outside it.\footnote{DHCP lease expiry faces the same structural
problem---``lease expired'' is not ``address unused''---but largely avoids it
through protocol design: the client contractually agrees to stop using the
address at expiry, and Duplicate Address Detection (RFC~5227) catches most
residual conflicts. The type error exists but the protocol mitigates it
bilaterally rather than ignoring it. Contrast Cassandra's
\texttt{gc\_grace\_seconds}, where no such mitigation exists---the wall-clock
constant is the entire mechanism.}

\section{Applying the Principles: CRDT Tombstone GC}

CRDT tombstone garbage collection is a clean test case for these principles,
because it forces every question they raise.\footnote{CRDT (Conflict-free
Replicated Data Type)---a family of data structures that can be replicated across
multiple nodes, updated independently and concurrently, and merged
deterministically without coordination. The merge function is mathematically
guaranteed to converge (it forms a join-semilattice), which eliminates the need
for consensus on every write---at the cost of carrying metadata, including
tombstones, that can only be safely discarded when all replicas have observed it.
See Shapiro et al.~\cite{shapiro2011comprehensive}.}

CRDTs track deletions with tombstones---markers that say ``this key was deleted
at time~$T$.'' A tombstone can only be GCed when \emph{every} replica has
observed the deletion, because removing it prematurely from any replica lets the
pre-delete value reappear (the zombie resurrection problem). ``Every replica has
observed it'' requires knowing two things: the group membership, and each
member's progress. The first is a semantic fact about the physical world. The
second is measurable.

\textbf{The slow-vs-dead problem is inescapable here.} A peer stops
acknowledging---is it partitioned, overloaded, or decommissioned? This is
exactly the FLP regime: no amount of observation from inside the system resolves
it. Any system that GCs tombstones automatically must promote ``unresponsive''
to ``gone''---an automatic semantic decision with unbounded consequences (silent
data resurrection if wrong).

\textbf{Cassandra's approach crosses every boundary.}
\texttt{gc\_grace\_seconds} is a wall-clock constant that assumes repair will
propagate tombstones within 10~days. This is a semantic decision disguised as a
configuration parameter, made automatically by the system with insufficient
information. The recovery action---hinted handoff replay on node return---violates
the recovery invariant by bursting accumulated writes into an already-stressed
cluster.

\textbf{What the principles demand instead:}

\begin{enumerate}
  \item \textbf{Separate what the system can measure from what it cannot.} The
    system \emph{can} measure: how many tombstones exist, how old they are,
    which peers have acknowledged which state, and how far behind each peer is.
    The system \emph{cannot} determine: whether an unresponsive peer will return.

  \item \textbf{Surface the measurable; defer the unmeasurable.} The API exposes
    what the system knows: \texttt{TombstoneCount()} reports pressure---how
    urgently GC is needed. \texttt{Status()} reports per-peer replication
    lag---who is behind and by how much. \texttt{BlockedBy(d)} names the specific
    peers preventing a given deletion from being GCed---the bottleneck, not just
    the symptom.

  \item \textbf{Reserve the semantic decision for the entity with sufficient
    information.} \texttt{Remove\-Peer()} is the only call that declares a node
    gone. It requires the operator---the one entity that can check whether the
    hardware is dead, whether the network is partitioned, whether the node is
    being decommissioned. The system will not make this call on its own.

  \item \textbf{Make the recovery action satisfy the invariant.} If a removed
    peer returns, it performs a full state transfer rather than delta replay. The
    returning node bears the cost of reconciliation, and that cost is bounded by
    the current state size---it does not scale with the duration of absence the
    way accumulated hints do. The cluster does not experience a burst.
\end{enumerate}

The result is a compensation boundary drawn where the information boundary
actually is: the system compensates for things it can measure (replication lag,
tombstone pressure) and surfaces everything else. No wall-clock constant
substitutes for operator judgment. No automatic action has unbounded
consequences.

\section{Design Rules}

\begin{enumerate}
  \item \textbf{Separate local decisions from semantic decisions.} TCP can kill a
    connection; only an operator can kill a node.

  \item \textbf{Derive timeouts from measured properties.} If you cannot measure
    it, you cannot timeout on it safely.

  \item \textbf{Check all three levels of the recovery invariant.} Individual
    cost bounded (level~1) is necessary but insufficient. Verify that aggregate
    recovery rate stays within freed capacity (level~2) and that recovery actions
    avoid amplifying the failure condition (level~3). Level~3 failure is a
    cascade seed.

  \item \textbf{Surface state, do not hide it.} When you cannot decide safely,
    expose what you know to whoever can.

  \item \textbf{Make failure a first-class API.} TCP does not hide
    retransmissions from the application---it surfaces \texttt{ETIMEDOUT}. Your
    system should not hide peer lag from the operator.

  \item \textbf{Prefer recoverable pessimism over unrecoverable optimism.} A
    slow operator is recoverable. Silent data loss is not.
\end{enumerate}

\section{Related Work}
\label{sec:related-work}

The phenomenon of recovery actions amplifying failures is well-identified across
multiple communities.

\paragraph{Metastable failures.} Bronson et al.~\cite{bronson2021metastable}
coined the term and identified the sustaining feedback loop---not the
trigger---as the root cause. Huang et al.~\cite{huang2022metastable} studied 22
metastable failures across 11 organizations, finding that ${>}50\%$ involved
retry storms. Isaacs and Alvaro~\cite{isaacs2025analyzing} propose a modeling
pipeline for discovering susceptibility at design time. The formal analysis by
Alvaro et al.~\cite{alvaro2025formal} provides a CTMC-based spectral
characterization of metastable states. This paper approaches the same phenomenon
from control theory rather than stochastic processes: the gain matrix asks
``does the perturbation converge or diverge?'' where CTMCs ask ``how long does
the system stay trapped?''

\paragraph{Resilience engineering.} David Woods' concept of
\emph{decompensation}---the exhaustion of adaptive capacity when recovery
mechanisms become the dominant source of
load~\cite{woods2015four,woods2018theory}---captures the qualitative phenomenon
that level~3 of the recovery invariant formalizes quantitatively. Woods asks
whether the system can extend capacity at all (\emph{graceful extensibility});
the compensation boundary asks \emph{where} the system draws the line between
reporting and acting.

\paragraph{Control theory for computing.} Hellerstein et
al.~\cite{hellerstein2004feedback} applied control theory (MIMO, transfer
functions, gain, pole placement) to computing systems. Their focus is on
designing controllers (admission control, scheduling), not on analyzing when
existing recovery mechanisms cascade. The gain matrix formalization in this paper
uses the same mathematical tradition for a different question.

\paragraph{Cascading failures in other domains.} Ramirez, Odijk, and
Bauso~\cite{ramirez2023cascading} applied gain matrix and Lyapunov stability
analysis to cascading failures in financial networks. The spectral radius
stability condition $\rho(G) < 1$ is the same; the domain and the specific
construction of $G$ differ.

\paragraph{Failure detectors.} Chandra and Toueg~\cite{chandra1996unreliable}
formalized the properties (completeness, accuracy) that failure detectors need to
enable consensus, establishing the theoretical foundation for the layered
decision authority model.

\paragraph{Exponential backoff stability.} The formal stability of binary
exponential backoff is well-studied. This paper uses TCP as a pedagogical
reference design---an existence proof that safe automatic recovery is
achievable---rather than making new claims about TCP itself.

\appendix

\section{The Gain Matrix: Formalizing the Feedback Bound}
\label{sec:gain-matrix}

Define the $n \times n$ recovery gain matrix $\mathbf{G}$:
\begin{equation}
  G_{ij} = \frac{\partial r_i}{\partial (-h_j)}
\end{equation}

Entry $G_{ij}$ measures how much additional recovery bandwidth appears on
resource $i$ when headroom on resource $j$ decreases by one unit. For systems
without backoff, $G_{ij} \geq 0$: less headroom means equal or greater recovery
demand.

When headroom drops by $\Delta\mathbf{h}$, recovery demand increases by
$\mathbf{G} \cdot \Delta\mathbf{h}$, reducing headroom further:
\begin{equation}
  \Delta\mathbf{h} \to \mathbf{G} \cdot \Delta\mathbf{h}
  \to \mathbf{G}^2 \cdot \Delta\mathbf{h} \to \cdots
\end{equation}
\begin{equation}
  \text{Total: } (\mathbf{I} + \mathbf{G} + \mathbf{G}^2 + \cdots)
  \cdot \Delta\mathbf{h}
  = (\mathbf{I} - \mathbf{G})^{-1} \cdot \Delta\mathbf{h}
\end{equation}

This geometric series converges if and only if $\rho(\mathbf{G}) < 1$, where
$\rho(\mathbf{G})$ is the spectral radius of $\mathbf{G}$ (the largest absolute
value among its eigenvalues).

\begin{quote}
\textbf{Stability condition:} $\rho(\mathbf{G}) < 1$.
\end{quote}

For a single resource, $\mathbf{G}$ is scalar and the condition reduces to
$G < 1$. For multiple resources, the spectral radius captures cross-resource
feedback: even if each diagonal entry $G_{ii} < 1$, the system cascades when
off-diagonal entries are large enough---$G_{12} \cdot G_{21} \geq 1$ means
resource~1's recovery stresses resource~2, whose recovery stresses resource~1.

\subsection{Monotonicity of the Stability Boundary}

For systems with $\mathbf{G} \geq 0$, the stability boundary is monotone in
load. This follows from a classical result.

\begin{theorem}[Perron-Frobenius, non-negative case]
Let $A \in \mathbb{R}^{n \times n}$ with $a_{ij} \geq 0$ for all $i, j$. Then:
\begin{enumerate}
  \item $\rho(A)$ is an eigenvalue of $A$.
  \item There exists $x \geq 0$, $x \neq 0$, with $Ax = \rho(A)x$.
  \item If $B \geq A$ entrywise ($b_{ij} \geq a_{ij}$ for all $i, j$), then
    $\rho(B) \geq \rho(A)$.
\end{enumerate}
\end{theorem}

\begin{proof}[Proof of property 3]
The Collatz-Wielandt characterization gives
$\rho(A) = \max_{x > 0} \min_i (Ax)_i / x_i$. If $B \geq A$ entrywise, then
$(Bx)_i \geq (Ax)_i$ for every $x > 0$; the inner minimum can only increase,
so the outer maximum can only increase.
\end{proof}

\paragraph{Consequence.} If load increase causes any $G_{ij}$ to grow,
$\rho(\mathbf{G})$ grows. For systems where $\mathbf{G}(\mathbf{h})$ is
continuous and recovery cost grows without bound as headroom vanishes---true for
every system in this paper except TCP---the intermediate value theorem gives a
critical headroom $\mathbf{h}^*$ where
$\rho(\mathbf{G}(\mathbf{h}^*)) = 1$:

\begin{itemize}
  \item $\mathbf{h}(t) > \mathbf{h}^*$: $\rho(\mathbf{G}) < 1$ $\to$ stable
  \item $\mathbf{h}(t) < \mathbf{h}^*$: $\rho(\mathbf{G}) > 1$ $\to$ cascade
\end{itemize}

The transition is one-directional. The system cannot self-rescue by becoming more
overloaded.

\section{Computing $h^*$: UNIX Swap Thrashing}
\label{sec:swap}

Swap crosses resource domains---it trades memory capacity (stock) for disk
bandwidth (flow)---making it a clean test of the multi-resource framework.

\begin{center}
\begin{tabular}{@{}lll@{}}
\toprule
Symbol & Meaning & Unit \\
\midrule
$S$ & Physical memory & pages \\
$W$ & Total working set (all processes) & pages \\
$E = W - S$ & Excess (pages that must swap) & pages \\
$f$ & Distinct page access rate (all processes) & pages/sec \\
$D'$ & Spare disk IOPS (total minus normal workload) & IOPS \\
\bottomrule
\end{tabular}
\end{center}

Under uniform access, the probability of hitting a swapped-out page is $E/W$.
Each fault costs 2~IOPS (swap-in read + swap-out write):
\begin{equation}
  P = \frac{2fE}{W}
\end{equation}

Setting $P = D'$ (disk saturation) and solving for the critical excess:
\begin{equation}
  E^* = \frac{D' W}{2f}
\end{equation}

For $W \approx S$, the swap margin ratio---excess as a fraction of physical
memory:
\begin{equation}
  \frac{E^*}{S} = \frac{D'}{2f}
\end{equation}

Computed for a 64\,GB server ($S = 16\text{M}$ pages at 4\,KB/page):

\begin{center}
\begin{tabular}{@{}llll@{}}
\toprule
Workload & $f$ (pages/sec) & HDD ($D' = 100$) & NVMe ($D' = 50{,}000$) \\
\midrule
Light web (1K req/s) & 10,000 & 320\,MB (0.5\%) & no $h^*$\dag \\
Database (1K qps) & 100,000 & 32\,MB (0.05\%) & 16\,GB (25\%) \\
Analytics scan & 1,000,000 & 3.2\,MB (0.005\%) & 1.6\,GB (2.5\%) \\
\bottomrule
\end{tabular}
\end{center}

\dag\ $E^*/S > 1$: disk bandwidth exceeds maximum possible fault rate. The
binding constraint shifts to latency (${\sim}100$\,\textmu s per NVMe fault),
which the bandwidth model does not capture.

On a database server with HDD, the system thrashes at 32\,MB of excess working
set---0.05\% over physical memory. On NVMe, 16\,GB. The 500$\times$ IOPS
improvement buys a proportional increase in margin, but the margin is finite and
unmonitored.

\paragraph{Gain matrix at $h^*$.} At disk saturation, the feedback loop:
\begin{equation}
  \mathbf{G} = \begin{pmatrix} \varepsilon & \alpha \\ \beta & \varepsilon
  \end{pmatrix}
  \quad\text{(rows: disk, memory)}
\end{equation}

\begin{itemize}[nosep]
  \item $\alpha = \partial(\text{swap IOPS}) / \partial(-h_{\mathrm{mem}})
    \approx 2f/W$: each page of excess generates $2f/W$ IOPS
  \item $\beta = \partial(\text{mem demand}) / \partial(-h_{\mathrm{disk}})$:
    disk saturation blocks processes, which hold pages plus ${\sim}8$--16\,KB of
    kernel I/O buffers per pending operation
\end{itemize}

Eigenvalues: $\varepsilon \pm \sqrt{\alpha\beta}$. Cascade when
$\alpha\beta \geq (1 - \varepsilon)^2$---the cross-resource loop (memory pressure
$\to$ swap $\to$ disk saturation $\to$ blocked processes $\to$ memory pressure)
has compound gain $\geq 1$.

\section{Computing $h^*$: TCP (The Degenerate Case)}
\label{sec:tcp-hstar}

TCP is the degenerate case: single resource (link bandwidth), scalar gain.

After $k$ consecutive timeouts, RTO doubles: $\mathrm{RTO}_k = \mathrm{RTO}_{\mathrm{base}} \times 2^k$.
Retransmission rate:
\begin{equation}
  r_k = \frac{1}{\mathrm{RTO}_{\mathrm{base}} \times 2^k}
\end{equation}

When congestion increases ($h$ decreases, another round of losses), $k$
increments:
\begin{equation}
  r_{k+1} = r_k / 2
\end{equation}

$G_{\mathrm{TCP}} = -0.5$. Each unit of congestion halves recovery demand.

Perron-Frobenius does not apply---$G$ has a negative entry.
$\rho(|G|) = 0.5 < 1$ at all loads. No $h^*$ exists. TCP is unconditionally
stable, not because recovery is bounded, but because it is a
\textbf{contraction}: recovery reduces the condition that triggered it.

\begin{center}
\begin{tabular}{@{}lll@{}}
\toprule
 & $G \geq 0$ (no backoff) & $G < 0$ (backoff) \\
\midrule
$\rho(G)$ monotone in load & Non-decreasing & Can decrease \\
$h^*$ exists & Always & May not \\
Self-rescue under overload & Impossible & Possible \\
\bottomrule
\end{tabular}
\end{center}

\section{Note on Nonlinearity}
\label{sec:nonlinearity}

The gain matrix $\mathbf{G}$ is the Jacobian of the recovery dynamics at the
current operating point---a local linearization. Three results bound the gap
between the linear model and the nonlinear system.

\paragraph{Local validity.} The Hartman-Grobman
theorem~\cite{hartman1960lemma,grobman1959homeomorphisms} guarantees that the
nonlinear system's local phase portrait is topologically equivalent to the
linearization when no eigenvalue of $\mathbf{G}$ lies on the unit circle.
$\rho(\mathbf{G}) < 1$ implies local asymptotic stability;
$\rho(\mathbf{G}) > 1$ implies local instability.

\paragraph{Monotonicity.} For $\mathbf{G} \geq 0$, $\rho(\mathbf{G})$ is
non-decreasing in every entry (Perron-Frobenius property~3). If recovery becomes
relatively more expensive as headroom shrinks---true for every system in this
paper except TCP---then $\rho(\mathbf{G})$ is non-decreasing in load. The
boundary $\mathbf{h}^*$ is a clean phase transition: stable on one side, cascade
on the other. The nonlinear system may reach instability sooner than the
linearization predicts, but it cannot reach stability later. The linear analysis
is a lower bound on stability---the safe direction for engineering.

\paragraph{Quantitative margins.} For exact convergence rates, overshoot, or
dynamics near $\rho(\mathbf{G}) = 1$, Lyapunov stability analysis provides the
general framework: construct $V(x)$ decreasing along trajectories; existence
proves stability, rate of decrease bounds
convergence~\cite{hirsch1985systems,smith1995monotone}. For the monotone case,
the theory of cooperative dynamical systems gives stronger structure: trajectories
are ordered, limit sets are equilibria, and the worst-case trajectory is
computable by evaluating $\mathbf{G}$ at its maximum entries.

\bibliographystyle{plain}
\bibliography{references}

\bigskip
\noindent\textit{Thanks for reading. If you work on distributed systems and have
seen the recovery invariant violated in ways not covered here, or if you spot
errors in the formalization, I'd like to hear about it.}

\end{document}
